\documentclass[interno]{gpt}

\titolo{Retrospettiva sprint 1}
\data{2026-11-18}
\orario{18:00}{20:00}
\luogo{Discord}
\redattore{spagnoli}
\verificatore{guglielmo}
%\assenti{}

\componenti{biasini,cescobolla,corro,guglielmo,mazzurana,spagnoli}
\begin{document}

\section{Ordine del giorno}
\begin{enumerate}
  \item Discussione dopo l'incontro con Eggon.
  \item Retrospettiva sprint 1.
  \item Definizione obiettivi e preventivo per sprint 2.
  \item Considerazioni sull'uso del diagramma di Gantt.
\end{enumerate}

\clearpage

\section{Svolgimento}

\subsection{Discussione dopo l'incontro con Eggon}
Dopo aver presentato al proponente i requisiti individuati e aver discusso dei vari dubbi emersi, il gruppo è in attesa di un feedback. I dettagli sull'incontro sono riportati nel verbale esterno \textit{VE-2026-11-18}.

\subsection{Retrospettiva sprint 1}
Il gruppo ha constatato che quasi tutti gli obiettivi dello sprint 1 sono stati raggiunti; viene posticipato allo sprint 2 soltanto il completamento del documento \textit{Analisi dei requisiti}, con l'aggiunta dei casi d'uso relativi ai requisiti individuati. Durante la retrospettiva sono inoltre emerse criticità legate alla rendicontazione oraria e alle stime iniziali, approfondite nella sezione 5.2.1.7 del \textit{Piano di progetto}.

\subsection{Definizione obiettivi e preventivo per sprint 2}
Il gruppo ha definito gli obiettivi da raggiungere per questo sprint, dando priorità al completamento del documento \textit{Analisi dei requisiti} e avviando la prima stesura del \textit{Piano di qualifica}. Per quanto riguarda il preventivo, ruoli e monte ore restano quelli pianificati, mantenendo le 14,75 ore non consumate nei budget dei rispettivi ruoli.

\subsection{Considerazioni sull'uso del diagramma di Gantt}
Il team ha valutato l'utilizzo del diagramma di Gantt per monitorare le attività del progetto e prevede di testarlo nel corso di questo sprint.

\section{Decisioni}
\begin{decisioni}
    \decisione{Modifica dei requisiti stilati in base al feedback del proponente.}
    \decisione{Produzione del documento \textit{Analisi dei requisiti}.}
    \decisione{Implementazione del diagramma di Gantt.}
    \decisione{Prima stesura del documento \textit{Piano di qualifica}.}
\end{decisioni}

\section{Azioni da intraprendere}
\begin{azioni}
  \azione[\#117]{Stesura verbale interno della riunione odierna.}{spagnoli}{19/11/2026}
  \azione[\#118]{Stesura verbale esterno dell'incontro odierno.}{biasini}{25/11/2026}
  \azione[\#119]{Produzione diario di bordo.}{biasini}{22/11/2026}
  \azione[\#121]{Revisione dei requisiti e produzione \textit{Analisi dei requisiti}.}{spagnoli, biasini, corro}{30/11/2026}
  \azione[\#122]{Prima stesura \textit{Piano di qualifica}.}{guglielmo, cescobolla, mazzurana}{30/11/2026}
  \azione[--]{Implementazione diagramma di Gantt.}{mazzurana}{30/11/2026}
\end{azioni}

\section{Prossima riunione}
La prossima riunione interna è prevista per lunedì 23 novembre 2026 alle ore 18:00, con lo scopo di discutere del feedback del proponente riguardo ai requisiti stilati. 

\end{document}